\bookStart{Genesis B}
\setBookCode{GenesisB}

\begin{flushright}%
\textbf{Dating:} C9th

\textbf{Meter:} \Fornyrdislag%para
\end{flushright}%

\section{Introduction}

\textbf{Genesis B} represents a part of the Old English Genesis translated from an Old Saxon original, presently represented by the fragmentary \textlink{SaxonGenesis} (edited above).  Most of the Saxon original underlying \emph{Genesis B} is now lost, but a few lines shared with \textlink{SaxonGenesis} attest to the primacy of the latter poem.

TODO: lack of alliteration in certain lines, notable Saxonisms.

\sectionline

\clearpage

\section{Text}

\bvg\bva[5][235]%
\llap{„}Ac níotað inc þæs \alst{ȯ}ðres ealles, \hld\ for·lǽtað þone \alst{ǽ}nne béam, &
\alst{w}ariað inc wið þone \alst{w}æstm. \hld\ Ne wyrð inc \alst{w}ilna gǽd.“ &
\alst{H}nigon þá mid \alst{h}eafdum \hld\ \alst{h}eofon-cyninge &
\alst{g}eorne to·\alst{g}éanes \hld\ and sǽdon ealles þanc, &
\alst{l}ista and þára \alst{l}ára. \hld\ Hé lét héo þæt \alst{l}and búan, &
\alst{h}wærf him þá tó \alst{h}eofenum \hld\ \alst{h}álig drihten, &
\alst{st}íð-ferhð cyning. \hld\ \alst{St}ód his hand-ge·weorc &
\alst{s}omod on \alst{s}ande, \hld\ nyston \alst{s}orga wiht &
to be·\alst{g}rornianne, \hld\ bútan héo \alst{g}odes willan &
\alst{l}engest \alst{l}ǽsten. \hld\ Héo wǽron \alst{l}éof gode &
þenden héo his \alst{h}álige word \hld\ \alst{h}ealdan woldon.\eva

\bvb TODO: Translation\evb\evg


\bvg\bva[6][246]%
Hæfde sé \alst{a}l-walda \hld\ \alst{e}ngel-cynna &
þurh \alst{h}and-mægen, \hld\ \alst{h}álig drihten, &
\alst{t}íene ge·\alst{t}rimede, \hld\ þǽm hé ge·\alst{t}rúwode wel &
þæt híe his giongor-scipe \hld\ fyligan wolden, &
\alst{w}yrcean his \alst{w}illan, \hld\ for·þon hé him ge·\alst{w}it for·geaf &
and mid his \alst{h}andum ge·sceop, \hld\ \alst{h}álig drihten. &
Ge·sett hæfde hé hie swá ge·sælig-lice, \hld\ ǽnne hæfde he swa swíðne ge·worhtne, &
swá mihtigne on his mod-ge·þóhte, \hld\ hé lét hine swá micles wealdan, &
héhstne tó him on heofona ríce, \hld\ hæfde hé hine swá hwítne ge·worhtne, &
swa wyn-lic wæs his wæstm on heofonum \hld\ þæt him cóm from weroda drihtne, &
ge·líc wæs hé þám leohtum steorrum. \hld\ Lof sceolde hé drihtnes wyrcean, &
dyran sceolde hé his dréamas on heofonum, \hld\ and sceolde his drihtne þancian &
þæs léanes þe hé him on þám leohte ge·scerede \hld\ þonne lǽte hé his hine lange wealdan. &
Ac hé wende hit him to wyrsan þinge, \hld\ on·gan him winn up hebban &
wið þone héhstan heofnes waldend, \hld\ þe siteð on þám hálgan stóle. &
Déore wæs hé drihtne úre; \hld\ ne mihte him be·dyrned weorðan &
þæt his engyl on·gan \hld\ ofer-mód wesan, &
á·hóf hine wið his herran, \hld\ sóhte hete-sprǽce, &
gylp-word on·gean, \hld\ nolde gode þeowian, &
cwæð þæt his líc wǽre \hld\ leoht and scéne, &
hwít and hiow-beorht. \hld\ Ne meahte hé æt his hige findan &
þæt gode wolde \hld\ geonger-dóme, &
þéodne þeowian. \hld\ Þu̇hte him sylfum &
þæt he mægyn and cræft \hld\ máran hæfde &
þonne sé hálga god \hld\ habban mihte &
folc-ge·stælna. \hld\ Feala worda ge·spæc &
se engel ofer-módes. \hld\ Þȯhte þurh his ánes cræft &
hú hé him streng-licran \hld\ stól ge·worhte, &
héahran on heofonum; \hld\ cwæð þæt hine his hige speone &
þæt hé west and norð \hld\ wyrcean on·gunne, &
trymede ge·timbro; \hld\ cwæð him tweo þu̇hte &
þæt hé gode wolde \hld\ geongra weorðan. &
\llap{„}Hwæt sceal ic winnan?“ cwæð he. \hld\ \llap{„}Nis me wihtæ þearf &
hearran to habbanne. \hld\ Ic mæg mid handum swá fela &
wundra ge·wyrcean. \hld\ Ic hæbbe ge·weald micel &
to gyrwanne \hld\ god-lecran stól, &
hearran on heofne. \hld\ Hwý sceal ic æfter his hyldo þeowian, &
búgan him swilces geongor-dómes? \hld\ Ic mæg wesan god swá hé. &
Big·standað mé strange ge·néatas, \hld\ þá ne willað mé æt þam stríðe ge·swícan, &
hæleþas heard-móde. \hld\ Hie habbað mé to hearran ge·corene, &
rófe rincas; \hld\ mid swilcum mæg man rǽd ge·þencean, &
fón mid swilcum folc-ge·steallan. \hld\ Frynd synd hie míne georne, &
holde on hyra hyge-sceaftum. \hld\ Ic mæg hýra hearra wesan, &
rǽdan on þis ríce. \hld\ Swá mé þæt riht ne þinceð, &
þæt ic ó·leccan \hld\ á·wiht þurfe &
gode æfter gode ǽnegum. \hld\ Ne wille ic leng his geongra wurþan.“ &
Þá hit sé all-walda \hld\ eall ge·hýrde, &
þæt his engyl on·gan \hld\ ofer·méde micel &
á·hebban wið his hearran \hld\ and spræc héa-lic word &
dol-líce wið drihten sinne, \hld\ sceolde hé þá dǽd on·gyldan, &
worc þæs ge·winnes ge·dǽlan, \hld\ and sceolde his wíte habban, &
ealra morðra mǽst. \hld\ Swa deð monna ge·hwilc &
þe wið his waldend \hld\ winnan on·gynneð &
mid máne wið þone mǽran drihten. \hld\ Þá wearð se mihtiga ge·bolgen, &
héhsta heofones waldend, \hld\ wearp hine of þán héan stóle. &
Hete hæfde hé æt his hearran ge·wunnen, \hld\ hyldo hæfde his fer·lorene, &
gram wearð him se góda on his móde. \hld\ For·þon hé sceolde grund ge·sécean &
heardes helle-wítes, \hld\ þæs þe hé wann wið heofnes waldend. &
Á·cwæð hine þá fram his hyldo \hld\ and hine on helle wearp, &
on þá déopan dala, \hld\ þær hé to déofle wearð, &
se féond mid his ge·férum eallum. \hld\ Feollon þá ufon of heofnum &
þurh-longe swá \hld\ þréo niht and dagas, &
þá englas of heofnum on helle, \hld\ and heo ealle for·scéop &
drihten to deoflum. · For·þon héo his dǽd and word &
noldon weorðian, \hld\ for·þon hé héo on wyrse leoht &
under eorðan neoðan, \hld\ æll-mihtig god, &
sette sige-léase \hld\ on þá sweartan helle. &
Þǽr hæbbað héo on ǽfyn \hld\ un-ge·met lange, &
ealra feonda ge·hwilc, \hld\ fýr ed-néowe, &
þonne cymð on úhtan \hld\ éasterne wind, &
forst fyrnum cald. \hld\ Symble fýr oððe gár, &
sum heard ge·wrinc \hld\ habban sceoldon. &
Worhte man hit him tó wíte, \hld\ ---hyra wor-uld wæs ge·hwyrfed---, &
forman sïðe, \hld\ fylde helle &
mid þám and·sacum. \hld\ Heoldon englas forð &
heofon-ríces héhðe, \hld\ þe ǽr godes hyldo ge·lǽston. &
Lágon þa ȯðre fynd on þám fýre, \hld\ þe ǽr swa feala hæfdon &
ge·winnes wið heora waldend. \hld\ Wíte þoliað, &
hátne heaðo-welm \hld\ helle to·middes,\eva

\bvb TODO: Translation\evb\evg


\bvg\bva[7][325]%
brand and bráde lígas, \hld\ swilce eac þá biteran récas, &
þrosm and þystro, \hld\ for·þon hie þegn-scipe &
godes for·gymdon. \hld\ Hie hyra gál be·swác — &
engles ofer·hygd — \hld\ noldon al-waldan &
word weorþian, \hld\ hæfdon wíte micel, &
wǽron þá be·feallene \hld\ fýre to botme &
on þá hátan hell \hld\ þurh hyge-léaste &
and þurh ofer·metto, \hld\ sóhton ȯþer land, &
þæt wæs léohtes léas \hld\ and wæs líges full, &
fýres fær micel. \hld\ Fynd on·géaton &
þæt híe hæfdon ge·wrixled \hld\ wíta un·rím &
þurh heora miclan mód \hld\ and þurh miht godes &
and þurh ofer·metto \hld\ ealra swíðost. &
Þá spræc sé ofer·móda cyning, \hld\ þe ǽr wæs engla scýnost, &
hwítost on heofne \hld\ and his hearran léof, &
drihtne dýre, \hld\ oð híe tó dole wurdon, &
þæt him for gál-scipe \hld\ god sylfa wearð &
mihtig on móde yrre. \hld\ Wearp hine on þæt morðer innan, &
niðer on þæt nío-bedd, \hld\ and scéop him naman siððan, &
cwæð sé héhsta \hld\ hátan sceolde &
Satan siððan \hld\ hét hine þǽre sweartan helle &
grundes gýman, \hld\ nalles wið god winnan. &
Satan maðelode, \hld\ sorgiende spræc, &
sé ðe helle forð \hld\ healdan sceolde, &
gíeman þæs grundes. \hld\ Wæs ǽr godes engel, &
hwít on heofne, \hld\ oð hine his hyge for·spéon &
and his ofer·metto \hld\ ealra swíðost, &
þæt hé ne wolde \hld\ wereda drihtnes &
word wurðian. \hld\ Weoll him on innan &
hyge ymb his heortan, \hld\ hát wæs him útan &
wráð-lic wíte. \hld\ Hé þá worde cwæð: &
\llap{„}Is þæs ǽnga styde \hld\ un·ge·líc swíðe &
þám ȯðrum þe we ǽr cu̇ðon, &
héan on heofon-ríce \hld\ þe me mín hearra on·lag, &
þeah wé hine for þám al-waldan \hld\ ágan ne móston, &
rómigan u̇res ríces. \hld\ Næfð hé þeah riht ge·dón &
þæt he u̇s hæfð be·fælled \hld\ fýre tó botme, &
helle þære hátan, \hld\ heofon-ríce be·numen; &
hafað hit ge·mearcod \hld\ mid mon-cynne &
to ge·settanne. \hld\ Þæt mé is sorga mǽst, &
þæt Adam sceal, \hld\ þe wæs of eorðan ge·worht, &
minne strong-lícan \hld\ stól be·healdan, &
wesan him on wynne, \hld\ and wé þis wíte þolien, &
hearm on þisse helle. \hld\ Wá lá, áhte ic mínra handa ge·weald &
and móste áne tíd \hld\ úte weorðan, &
wesan áne winter-stunde, \hld\ þonne ic mid þys werode--- &
Ac licgað me ymbe \hld\ íren-benda, &
rídeð racentan sál. \hld\ Ic eom ríces léas; &
habbað mé swá hearde \hld\ helle clommas &
fæste be·fangen. \hld\ Hér is fýr micel, &
ufan and neoðone. \hld\ Ic á ne ge·seah &
láðran land-scipe. \hld\ Líg ne á·swámað, &
hát ofer helle. \hld\ Me habbað hringa ge·spong, &
slíð-hearda sál \hld\ sïðes á·myrred, &
á·fyrred mé min féðe; \hld\ fét synt ge·bundene, &
handa ge·hæfte. \hld\ Synt þissa heldora &
wegas for·worhte, \hld\ swá ic mid wihte ne mæg &
of þissum lioðo-bendum. \hld\ Licgað me ymbe &
heardes írenes \hld\ háte ge·slægene &
grindlas gréate. \hld\ Mid þý mé god hafað &
ge·hæfted be þám healse, \hld\ swá ic wát hé minne hige cu̇ðe; &
and þæt wiste éac \hld\ weroda drihten, &
þæt sceolde unc Adam \hld\ yfele ge·wurðan &
ymb þæt heofon-ríce, \hld\ þær ic áhte mínra handa ge·weald.\eva

\bvb TODO: Translation\evb\evg



\bvg\bva[8][389]%
Ac þoliað wé nu þréa on helle, \hld\ ---þæt syndon þystro and hæto---, &
grimme, grund-lease. \hld\ Hafað us god sylfa &
for·swapen on þás sweartan mistas; \hld\ swá hé us ne mæg ǽnige synne ge·stælan, &
þæt wé him on þám lande láð ge·fremedon, \hld\ he hæfð us þeah þæs leohtes be·scyrede, &
be·worpen on ealra wíta mǽste. \hld\ Ne magon we þæs wrace ge·fremman, &
ge·léanian him mid láðes wihte \hld\ þæt he us hafað þæs leohtes be·scyrede. &
He hæfð nú gem·earcod anne middan-geard, \hld\ þǽr he hæfð mon ge·worhtne &
æfter his on·líc-nesse. \hld\ Mid þám hé wile eft ge·settan &
heofona ríce mid hluttrum saulum. \hld\ Wé þæs sculon hycgan georne, &
þæt wé on Adame, \hld\ gif wé ǽfre mægen, &
and on his eafrum swá some, \hld\ andan ge·bétan, &
on·wendan him þǽr willan sínes, \hld\ gif wé hit mægen wihte á·þencan. &
Ne ge·lýfe ic mé nú þæs leohtes furðor \hld\ þæs þe he him þenceð lange niotan, &
þæs éades mid his engla cræfte. \hld\ Ne mágon wé þæt on aldre ge·winnan, &
þæt we mihtiges godes mód on·wæcen. \hld\ Uton oð·wendan hit nú monna bearnum, &
þæt heofon-ríce, nú wé hit habban ne móton, \hld\ ge·dón þæt híe his hyldo for·lǽten, &
þæt híe þæt on·wendon þæt hé mid his worde be·béad. \hld\ þonne weorð hé him wráð on móde, &
á·hwet híe from his hyldo. \hld\ þonne sculon híe þás helle sécan &
and þás grimman grundas. \hld\ Þonne móton wé híe u̇s to giongrum habban, &
fíra bearn on þissum fæstum clomme. \hld\ On·ginnað nú ymb þá fyrde þencean. &
Gif ic ǽnegum þegne \hld\ þéoden-mádmas &
géara for·géafe, \hld\ þenden we on þan godan rice &
gesælige sæton \hld\ and hæfdon u̇re setla ge·weald, &
þonne hé mé ná on léofran tíd \hld\ léanum ne meahte &
míne gife gyldan, \hld\ gif his gien wolde &
mínra þegna hwilc \hld\ ge·þafa wurðan, &
þæt hé up heonon \hld\ úte mihte &
cuman þurh þás clustro, \hld\ and hæfde cræft mid him &
þæt hé mid feðer-homan \hld\ fléogan meahte, &
windan on wolcne, \hld\ þǽr ge·worht stondað &
Adam and Eue \hld\ on eorð-ríce &
mid welan be·wunden, \hld\ and wé synd á·worpene hider &
on þás déopan dalo. \hld\ Nú híe drihtne synt &
wurðran micle, \hld\ and móton him þone welan a´gan &
þe wé on heofon-ríce \hld\ habban sceoldon, &
ríce mid rihte; \hld\ is sé rǽd ge·scyred &
monna cynne. \hld\ Þæt mé is on mínum mode swá sár, &
on mínum hyge hreoweð, \hld\ þæt híe heofon-ríce &
ágan tó aldre. \hld\ Gif hit eower ǽnig mæge &
ge·wendan mid wihte \hld\ þæt híe word godes &
láre for·lǽten, \hld\ sóna híe him þe láðran beoð. &
Gif híe brecað his ge·bod-scipe, \hld\ þonne hé him á·bolgen wurðeþ; &
siððan bið him sé wela on·wended \hld\ and wyrð him wíte ge·garwod, &
sum heard hearm-scearu. \hld\ Hycgað his ealle, &
hú gé hí be·swicen. \hld\ Siððan ic mé sefte mæg &
restan on þyssum racentum, \hld\ gif him þæt ríce losað. &
Sé þe þæt ge·lǽsteð, \hld\ him bið léan gearo &
æfter tó aldre, \hld\ þæs wé hér inne mágon &
on þyssum fýre forð \hld\ fremena gewinnan. &
Sittan lǽte ic hine wið mé sylfne, \hld\ swá hwá swá þæt secgan cymeð &
on þas hátan helle, \hld\ þæt hie heofon-cyninges &
un·wurð-líce \hld\ wordum and dǽdum &
láre \edtext{[...]}{\Afootnote{At this point two folios are missing.  These would have contained sections 9--10 of the poem.}}\eva

\bvb TODO: Translation\evb\evg


\bvg\bva[11][442]%
An·gan hine þa gyrwan \hld\ godes andsaca, &
fus on frætwum, \hld\ hæfde fæcne hyge, &
hæleðhelm on heafod asette \hld\ and þone ful hearde geband, &
spenn mid spangum; \hld\ wiste him spræca fela, &
wora worda. \hld\ Wand him up þanon, &
hwearf him þurh þa helldora, \hld\ ---hæfde hyge strangne---, &
leolc on lyfte \hld\ laþwendemod, &
swang þæt fyr on twa \hld\ feondes cræfte; &
wolde dearnunga \hld\ drihtnes geongran, &
mid man-dǽdum \hld\ men be·swícan, &
for·lǽdan and for·lǽran, \hld\ þæt híe wurdon láð gode. &
Hé þá ge·férde \hld\ þurh féondes cræft &
oððæt hé Adam \hld\ on eorð-ríce, &
godes hand-ge·sceaft, \hld\ gearone funde, &
wis-líce ge·worht, \hld\ and his wíf somed, &
fréo fægroste, \hld\ swá híe fela cu̇ðon &
godes ge·gearwigean, \hld\ þá him tó gingran self &
metot man-cynnes \hld\ mearcode selfa. &
And him bí twégin \hld\ béamas stódon &
þá wǽron útan \hld\ ofætes ge·hlædene, &
ge·wered mid wæstme, \hld\ swá híe waldend god, &
héah heofon-cyning \hld\ handum ge·sette, &
þæt þǽr yldo bearn \hld\ móste on céosan &
gódes and yfeles, \hld\ gumena æg-hwilc, &
welan and wáwan. \hld\ Næs sé wæstm ge·líc. &
Ȯðer wæs swá wyn-lic, \hld\ wlitig and scíene, &
liðe and lof-sum, \hld\ þæt wæs lífes béam; &
moste on ec-nisse \hld\ æfter lybban, &
wesan on wor-ulde, \hld\ sé þæs wæstmes on·bát, &
swá him æfter þý \hld\ yldo ne derede, &
ne suht swáre, \hld\ ac móste symle wesan &
lungre on lustum \hld\ and his líf ágon, &
hyldo heofon-cyninges \hld\ hér on wor-ulde, &
habban him tó wǽron \hld\ witot ge·þinge &
on þone héan heofon, \hld\ þonne hé* heonon wende. &
Þonne wæs sé ȯðer \hld\ eallenga sweart, &
dim and þystre; \hld\ þæt wæs déaðes béam, &
sé bær bitres fela. \hld\ Sceolde bú witan &
ylda æg-hwilc \hld\ yfles and gódes &
ge·wand on þisse worulde. \hld\ Sceolde on wíte á &
mid swáte and mid sorgum \hld\ siððan libban, &
swá hwá swá ge·byrgde \hld\ þæs on þám béame ge·weox. &
Sceolde hine yldo be·niman \hld\ ellen-dǽda, &
dréamas and driht-scipes, \hld\ and him beon déað scyred. &
Lýtle hwíle \hld\ sceolde hé his lífes níotan, &
sécan þonne landa \hld\ sweartost on fýre. &
Sceolde féondum þéowian, \hld\ þǽr is ealra frécna mǽste &
léodum tó langre hwíle. \hld\ Þæt wiste sé láða georne, &
dyrne deofles boda \hld\ þe wið drihten wann. &
Wearp hine þá on wyrmes líc \hld\ and wand him þá ymb-utan &
þone déaðes béam \hld\ þurh deofles cræft, &
ge·nam þǽr þæs ofætes \hld\ and wende hine eft þanon &
þǽr hé wiste hand-ge·weorc \hld\ heofon-cyninges. &
On·gon hine þá frínan \hld\ forman worde &
se láða mid lígenum: \hld\ \llap{„}Langað þé á·wuht, &
Adam, up tó gode? \hld\ Ic eom on his ǽrende hider &
feorran ge·féred, \hld\ ne þæt nú fyrn ne wæs &
þæt ic wið hine sylfne sæt. \hld\ Þá hét hé mé on þysne sïð faran, &
hét þæt þú þisses ofætes æte, \hld\ cwæð þæt þín abal and cræft &
and þín mód-sefa \hld\ mára wurde, &
and þín líc-homa \hld\ leohtra micle, &
þín ge·sceapu scénran, \hld\ cwæð þæt þe æniges sceates ðearf &
ne wurde on worulde. \hld\ Nu þu willan hæfst, &
hyldo geworhte \hld\ heofoncyninges, &
to þance geþenod \hld\ þínum hearan, &
hæfst þe wið drihten dyrne geworhtne. \hld\ Ic gehyrde hine þine dæd and word &
lofian on his leohte \hld\ and ymb þin lif sprecan. &
Swa þu læstan scealt \hld\ þæt on þis land hider &
his bodan bringað. \hld\ Brade synd on worulde &
grene geardas, \hld\ and god siteð &
on þam hehstan \hld\ heofna rice, &
ufan alwalda. \hld\ Nele þa earfeðu &
sylfa habban \hld\ þæt he on þysne sið fare, &
gumena drihten, \hld\ ac he his gingran sent &
to þinre spræce. \hld\ Nu he þe mid spellum het &
listas læran. \hld\ Læste þu georne &
his ambyhto, \hld\ nim þe þis ofæt on hand, &
bit his and byrge. \hld\ Þe weorð on þínum breostum rum, &
wæstm þy wlitegra. \hld\ Þe sende waldend god, &
þin hearra þas helpe \hld\ of heofon-ríce.“ &
Adam maðelode \hld\ þær he on eorðan stod, &
self-sceafte guma: \hld\ \llap{„}Þonne ic sige-drihten, &
mihtigne god, \hld\ mæðlan ge·hýrde &
strangre stemne, \hld\ and me her stondan het, &
his be·bodu healdan, \hld\ and me þas brýd forgeaf, &
wlite-scíene wíf, \hld\ and me warnian het &
þæt ic on þone déaðes béam \hld\ be·droren ne wurde, &
be·swicen tó swíðe, \hld\ he cwæð þæt þá sweartan helle &
healdan sceolde \hld\ sé ðe bí his heortan wuht &
láðes ge·lǽde. \hld\ Nát þeah þú mid lígenum fare &
þurh dýrne ge·þanc \hld\ þe þu drihtnes eart &
boda of heofnum. \hld\ Hwæt, ic þinra bysna ne mæg, &
worda ne wisna \hld\ wuht oncnawan, &
siðes ne sagona. \hld\ Ic wat hwæt he me self bebead, &
nergend user, \hld\ þa ic hine nehst geseah; &
he het me his word weorðian \hld\ and wel healdan, &
læstan his lare. \hld\ Þu gelic ne bist &
ænegum his engla \hld\ þe ic ær geseah, &
ne þu me oðiewdest \hld\ ænig tacen &
þe he me þurh treowe \hld\ to onsende, &
min hearra þurh hyldo. \hld\ Þy ic þe hyran ne cann. &
Ac þu meaht þe forð faran. \hld\ Ic hæbbe me fæstne geleafan &
up to þam ælmihtegan gode \hld\ þe me mid his earmum worhte, &
her mid handum sínum. \hld\ He mæg ﹒ me of his hean rice &
geofian mid goda gehwilcum, \hld\ þeah he his gingran ne sende.“\eva

\bvb TODO: Translation\evb\evg


\bvg\bva[12][547]%
Wende hine wráð-mód \hld\ þǽr hé þæt wíf ge·séah &
on eorð-ríce \hld\ Euan stondan, &
scéone ge·sceapene, \hld\ cwæð þæt sceaðena mǽst &
eallum heora eaforum \hld\ æfter siððan &
wurde on wor-ulde: \hld\ \llap{„}Ic wát, inc waldend god &
a·bolgen wyrð, \hld\ swá ic him þisne bod-scipe &
selfa secge, \hld\ þonne ic of þys sïðe cume &
ofer langne weg, \hld\ þæt git ne læstan wel &
hwilc ærende swa \hld\ he easten hider &
on þysne sið sendeð. \hld\ Nu sceal he sylf faran &
to incre andsware; \hld\ ne mæg his ærende &
his boda beodan; \hld\ þy ic wat þæt he inc abolgen wyrð, &
mihtig on mode. \hld\ Gif þu þeah mínum wilt, &
wíf willende, \hld\ wordum hýran, &
þu meaht his þonne rúme \hld\ rǽd ge·þencan. &
Ge·hyge on þínum breostum \hld\ þæt þu inc bam twam meaht &
wíte be·warigan, \hld\ swa ic þe wisie. &
Æt þisses ofetes; \hld\ þonne wurðað þin eagan swa leoht &
þæt þu meaht swa wide \hld\ ofer woruld ealle &
geseon siððan, \hld\ and selfes stol &
herran þines, \hld\ and habban his hyldo forð. &
Meaht þu Adame \hld\ eft gestyran, &
gif þu his willan hæfst \hld\ and he þínum wordum getrywð. &
Gif þu him to soðe sægst \hld\ hwylce þu selfa hæfst &
bisne on breostum, \hld\ þæs þu gebod godes &
lare læstes, \hld\ he þone laðan strið, &
yfel and·wyrde \hld\ an for·læteð &
on breost-cofan, \hld\ swa wit him bu tu &
an sped sprecað. \hld\ Span þu hine georne &
þæt he þine lare læste, \hld\ þy læs gyt lað gode, &
incrum waldende, \hld\ weorðan þyrfen. &
Gif þu þæt angin fremest, \hld\ idesa seo betste, &
forhele ic incrum herran \hld\ þæt me hearmes swa fela &
Adam gespræc, \hld\ eargra worda. &
Tyhð me untryowða, \hld\ cwyð þæt ic seo teonum georn, &
gramum ambyhtsecg, \hld\ nales godes engel. &
Ac ic cann ealle swa geare \hld\ engla gebyrdo, &
heah heofona gehlidu; \hld\ wæs seo hwil þæs lang &
þæt ic geornlice \hld\ gode þegnode &
þurh holdne hyge, \hld\ herran mínum, &
drihtne selfum; \hld\ ne eom ic déofle ge·líc.“ &
Lædde hie swa mid ligenum \hld\ and mid listum speon &
idese on þæt unriht, \hld\ oðþæt hire on innan ongan &
weallan wyrmes geþeaht, \hld\ ---hæfde hire wacran hige &
metod ge·mearcod---, \hld\ þæt heo hire mod ongan &
lǽtan æfter þam lárum; \hld\ forþon heo æt þam laðan onfeng &
ofer drihtnes word \hld\ deaðes beames &
weorcsumne wæstm. \hld\ Ne wearð wyrse dæd &
monnum gemearcod. \hld\ Þæt is micel wundor &
þæt hit ece god \hld\ æfre wolde &
þeoden þolian, \hld\ þæt wurde þegn swa monig &
forlædd be þam lygenum \hld\ þe for þam larum com. &
Heo þa þæs ofætes æt, \hld\ alwaldan bræc &
word and willan. \hld\ Þa meahte heo wide geseon &
þurh þæs laðan læn \hld\ þe hie mid ligenum beswac, &
dearnenga bedrog, \hld\ þe hire for his dædum com, &
þæt hire þuhte hwitre \hld\ heofon and eorðe, &
and eall þeos woruld wlitigre, \hld\ and geweorc godes &
micel and mihtig, \hld\ þeah heo hit þurh monnes geþeaht &
ne sceawode. \hld\ Ac se sceaða georne &
swicode ymb þa sawle \hld\ þe hire ær þa siene onlah, &
þæt heo swa wide \hld\ wlitan meahte &
ofer heofonrice. \hld\ Þa se forhatena spræc &
þurh feondscipe \hld\ ---nalles he hie freme lærde---: &
\llap{„}Þu meaht nu þe self geseon, \hld\ swa ic hit þe secgan ne þearf, &
Eue seo gode, \hld\ þæt þe is ungelic &
wlite and wæstmas, \hld\ siððan þu mínum wordum getruwodest, &
læstes mine lare. \hld\ Nu scineð þe leoht fore &
glæd-lic on·gean \hld\ þæt ic from gode brohte &
hwit of heofonum; \hld\ nu þu his hrinan meaht. &
Sæge Adame \hld\ hwilce þu gesihðe hæfst &
þurh minne cime cræfta. \hld\ Gif giet þurh cuscne siodo &
læst mina lara, \hld\ þonne gife ic him þæs leohtes genog &
þæs ic þe swa godes \hld\ gegired hæbbe. &
Ne wite ic him þa womcwidas, \hld\ þeah he his wyrðe ne sie &
to alætanne; \hld\ þæs fela he me laðes spræc.“ &
Swá hire eaforan sculon \hld\ æfter lybban: &
þonne híe láð ge·dóð, \hld\ híe sculon lufe wyrcean, &
bétan heora hearran hearm-cwide \hld\ ond habban his hyldo forð. &
Þá gíen tó Adame \hld\ idesa scíenost, &
wífa wlitegost \hld\ þe on wor-uld come, &
for·þon héo wæs hand-ge·weorc \hld\ heofon-cyninges, &
þeah héo þá dearnenga \hld\ fordon wurde, &
for·lǽd mid lígenum, \hld\ þæt híe láð gode &
þurh þæs wráðan ge·þanc \hld\ weorðan sceoldon, &
þurh þæs déofles searo \hld\ dóm for·lǽtan, &
hierran hyldo, \hld\ hefon-ríces þolian &
monige hwíle. \hld\ Bið þam men full wá &
þe hine ne warnað \hld\ þonne hé his ge·weald hafað. &
Sum héo hire on handum bær, \hld\ sum hire æt heortan læg, &
æppel un·sælga, \hld\ þone hire ǽr for·bead &
drihtna drihten, \hld\ déað-béames ofet, &
and þæt word á·cwæð \hld\ wuldres aldor, &
þæt þæt micle morð \hld\ menn ne þorfton &
þegnas þolian, \hld\ ac hé þéoda ge·hwám &
hefon-ríce for·geaf, \hld\ hálig drihten, &
wíd-brádne welan, \hld\ gif híe þone wæstm an &
lǽtan wolden \hld\ þe þæt láð-tréow &
on his bógum bær, \hld\ bitre ge·fylled; &
þæt wæs déaðes béam \hld\ þe him drihten for·bead. &
For·léc hie þá mid lígenum \hld\ sé wæs láð gode, &
on hete heofon-cyninges, \hld\ and hyge Euan, &
wífes wác ge·þoht, \hld\ þæt héo on·gan his wordum trúwian, &
lǽstan his láre, \hld\ and ge·léafan nom &
þæt hé þá bysene from gode \hld\ brungen hæfde &
þe hé hire swá wǽr-lice \hld\ wordum sægde, &
iewde hire tácen \hld\ and tréowa ge·hét, &
his holdne hyge. \hld\ Þá héo tó hire hearran spræc: &
\llap{„}Adam, fréa mín, \hld\ þis ofet is swá swéte, &
blíð on bréostum, \hld\ and þes boda scíene, &
godes engel gód, \hld\ ic on his gearwan ge·séo &
þæt hé is ǽrend-secg \hld\ uncres hearran, &
hefon-cyninges. \hld\ His hyldo is unc betere &
to ge·winnanne \hld\ þonne his wiðer·médo. &
Gif þú him héo·dæg wuht \hld\ hearmes ge·sprǽce, &
hé for·gifð hit þéah, \hld\ gif wit him geongor-dóm &
lǽstan willað. \hld\ Hwæt scal þé swá láð-lic stríð &
wið þínes hearran bodan? \hld\ Unc is his hyldo þearf; &
hé mæg unc ǽrendian \hld\ tó þám al-waldan, &
heofon-cyninge. \hld\ Ic mæg heonon ge·séon &
hwǽr hé sylf siteð, \hld\ ---þæt is su̇ð and éast---, &
welan be·wunden, \hld\ sé ðás wor-uld ge·scéop; &
ge·séo ic him his englas \hld\ ymbe hweorfan &
mid feðer-haman, \hld\ ealra folca mǽst, &
wereda wynsumast. \hld\ Hwá meahte me swelc ge·wit gifan, &
gif hit gegnunga \hld\ god ne on·sende, &
heofones waldend? \hld\ Gehyran mæg ic rume &
and swa wide ge·seon \hld\ on woruld ealle &
ofer þas sidan ge·sceaft, \hld\ ic mæg swegles gamen &
gehyran on heofnum. \hld\ Wearð me on hige leohte &
utan and innan, \hld\ siðþan ic þæs ofætes onbat. &
Nu hæbbe ic his her on handa, \hld\ herra se goda; &
gife ic hit þe georne. \hld\ Ic gelyfe þæt hit from gode come, &
broht from his bysene, \hld\ þæs me þes boda sægde &
wǽrum wordum. \hld\ Hit nis wuhte gelic &
elles on eorðan, \hld\ buton swa þes ar sægeð, &
þæt hit gegnunga \hld\ from gode come.“\eva

\bvb TODO: Translation\evb\evg



\bvg\bva[13][684]%
Hío spræc him þicce tó \hld\ and speon hine ealne dæg &
on þa dimman dæd \hld\ þæt hie drihtnes heora &
willan bræcon. \hld\ Stod se wraða boda, &
legde him lustas on \hld\ and mid listum speon, &
fylgde him frecne; \hld\ wæs se feond full neah &
þe on þa frecnan fyrd \hld\ gefaren hæfde &
ofer langne weg; \hld\ leode hogode &
on þæt micle morð \hld\ men forweorpan, &
forlæran and forlædan, \hld\ þæt hie læn godes, &
ælmihtiges gife \hld\ an forleten, &
heofenrices ge·weald. \hld\ Hwæt, se hellsceaða &
gearwe wiste \hld\ þæt hie godes yrre &
habban sceoldon \hld\ and hellgeþwing, &
þone nearwan nið \hld\ niede onfon, &
siððan hie gebod godes \hld\ forbrocen hæfdon, &
þa he forlærde \hld\ mid ligenwordum &
to þam unræde \hld\ idese sciene, &
wifa wlitegost, \hld\ þæt heo on his willan spræc, &
wæs hire on helpe \hld\ handweorc godes &
to forlæranne. &
Heo spræc ða to Adame \hld\ idesa sceonost &
ful þiclice, \hld\ oð þam þegne ongan &
his hige hweorfan, \hld\ þæt he þam gehate getruwode &
þe him þæt wif \hld\ wordum sægde. &
Heo dyde hit þeah þurh holdne hyge, \hld\ nyste þæt þær hearma swa fela, &
fyren-earfeða, \hld\ fylgean sceolde &
monna cynne \hld\ þæs héo on mód ge·nam &
þæt héo þæs láþan bodan \hld\ lárum hýrde &
ac wénde þæt héo hyldo \hld\ heofon-cyninges &
worhte mid þám wordum \hld\ þe héo þám were swelce &
tácen ó·þiewde \hld\ and treowe ge·hét  &
oð·þæt Adame \hld\ innan breostum &
his hyge hwyrfde \hld\ and his heorte on·gann &
wendan to hire willan. \hld\ He æt þam wife onfeng &
helle and hinnsið, \hld\ þeah hit nære haten swa, &
ac hit ofetes noman \hld\ agan sceolde; &
hit wæs þeah deaðes swefn \hld\ and deofles gespon, &
hell and hinnsið \hld\ and hæleða forlor, &
menniscra morð, \hld\ þæt hie to mete dædon, &
ofet unfæle. \hld\ Swa hit him on innan com, &
hran æt heortan, \hld\ hloh þa and plegode &
boda bitre gehugod, \hld\ sægde begra þanc &
hearran sínum: \hld\ \llap{„}Nu hæbbe ic þine hyldo me &
witode geworhte, \hld\ and þinne willan gelæst &
to ful monegum dæge. \hld\ Men synt forlædde, &
Adam and Eue. \hld\ Him is unhyldo &
waldendes witod, \hld\ nu hie wordcwyde his, &
lare forleton. \hld\ Forþon hie leng ne magon &
healdan heofonrice. \hld\ Ac hie to helle sculon &
on þone sweartan sið. \hld\ Swa þu his sorge ne þearft &
beran on þínum breostum, \hld\ þær þu gebunden ligst, &
murnan on mode, \hld\ þæt her men bun &
þone hean heofon, \hld\ þeah wit hearmas nu, &
þrea-weorc þoliað, \hld\ and þystre land, &
and þurh þin micle mod \hld\ monig for·léton &
on heofon-ríce \hld\ heah-ge·timbro, &
god-líce geardas. \hld\ Unc wearð god yrre &
for·þon wit him noldon \hld\ on heofon-ríce &
hnígan mid heafdum \hld\ hálgum drihtne &
þurh geongor-dóm; \hld\ ac unc ge·genge ne wæs &
þæt wit him on þegn-scipe \hld\ þeowian wolden. &
Forþon unc waldend wearð \hld\ wráð on móde, &
on hyge hearde, \hld\ and u̇s on helle be·dráf, &
on þæt fýr fylde \hld\ folca mǽste, &
and mid handum his \hld\ eft on heofon-ríce &
rihte rodor-stólas \hld\ and þæt ríce for·geaf &
monna cynne. \hld\ Mæg þín mód wesan &
blíðe on bréostum, \hld\ for·þon hér synt bú tú ge·dón: &
gé þæt hæleða bearn \hld\ heofon-ríces sculon &
leode for·lǽtan \hld\ and on þæt líg to þe &
háte hweorfan, \hld\ eac is hearm gode, &
mód-sorg ge·macod. \hld\ Swá hwæt swá wit hér morðres þoliað, &
hit is nú Adame \hld\ eall for·golden &
mid hearran hete \hld\ and mid hæleða for·lore, &
monnum mid morðes cwealme. \hld\ For·þon is min mód ge·hǽled, &
hyge ymb heortan ge·rúme, \hld\ ealle synt uncre hearmas ge·wrecene &
láðes þæt wit lange þoledon. \hld\ Nú wille ic eft þam líge néar, &
Satan ic þǽr sécan wille; \hld\ hé is on þǽre sweartan helle &
hæft mid hringa ge·sponne.“ \hld\ Hwearf him eft niðer &
boda bitresta; \hld\ sceolde he þa brádan lígas &
sécan helle ge·hliðo, \hld\ þær his hearra læg &
simon ge·sǽled. \hld\ Sorgedon ba twa, &
Adam and Eue, \hld\ and him oft betuh &
gnornword gengdon; \hld\ godes him ondredon, &
heora herran hete, \hld\ heofoncyninges nið &
swiðe onsæton; \hld\ selfe forstodon &
his word onwended. \hld\ Þæt wif gnornode, &
hof hreowigmod, \hld\ ---hæfde hyldo godes, &
lare forlæten---, \hld\ þa heo þæt leoht geseah &
ellor scriðan \hld\ þæt hire þurh untreowa &
tacen iewde \hld\ se him þone teonan geræd, &
þæt hie helle níð \hld\ habban sceoldon, &
hynða unrim; \hld\ forþam him higesorga &
burnon on breostum. \hld\ Hwilum to gebede feollon &
sinhiwan somed, \hld\ and sigedrihten &
gódne grétton \hld\ and god nemdon, &
heofones waldend, \hld\ and hine bædon &
þæt hie his hearmsceare \hld\ habban mosten, &
georne fulgangan, \hld\ þa hie godes hæfdon &
bodscipe abrocen. \hld\ Bare hie gesawon &
heora lichaman; \hld\ næfdon on þam lande þa giet &
sælða ge·setena, \hld\ ne hie sorge wiht &
weorces wiston, \hld\ ac hie wel meahton &
libban on þam lande, \hld\ gif hie wolden lare godes &
forweard fremman. \hld\ þa hie fela sprǽcon &
sorh-worda somed, \hld\ sin-hiwan twa. &
Adam ge·mǽlde \hld\ and to Euan spræc: &
\llap{„}Hwæt, þu Eue, hæfst \hld\ yfele gemearcod &
uncer sylfra sið. \hld\ Gesyhst þu nu þa sweartan helle &
grædige and gifre. \hld\ Nu þu hie grimman meaht &
heonane gehyran. \hld\ Nis heofonrice &
gelic þam lige, \hld\ ac þis is landa betst, &
þæt wit þurh uncres hearran þanc \hld\ habban moston. &
þær þu þam ne hierde \hld\ þe unc þisne hearm geræd, &
þæt wit waldendes \hld\ word forbræcon, &
heofoncyninges. \hld\ Nu wit hreowige magon &
sorgian for þis siðe. \hld\ Forþon he unc self bebead &
þæt wit unc wite \hld\ warian sceolden, &
hearma mǽstne. \hld\ Nu slit me hunger and þurst &
bitre on breostum, \hld\ þæs wit begra ær &
wæron orsorge \hld\ on ealle tid. &
Hú sculon wit nu libban \hld\ oððe on þis lande wesan, &
gif her wind cymð, \hld\ westan oððe eastan, &
suðan oððe norðan? \hld\ Gesweorc up færeð, &
cymeð hægles scur \hld\ hefone getenge, &
færeð forst on gemang, \hld\ se byð fyrnum ceald. &
Hwilum of heofnum \hld\ hate scineð, &
blicð þeos beorhte sunne, \hld\ and wit her baru standað, &
unwered wædo. \hld\ Nys unc wuht beforan &
to scursceade, \hld\ ne sceattes wiht &
to mete gemearcod, \hld\ ac unc is mihtig god, &
waldend wraðmod. \hld\ To hwon sculon wit weorðan nu? &
Nu me mæg hreowan \hld\ þæt ic bæd heofnes god, &
waldend þone godan, \hld\ þæt he þe her worhte to me &
of liðum mínum, \hld\ nu þu me forlæred hæfst &
on mines herran hete. \hld\ Swa me nu hreowan mæg &
æfre to aldre \hld\ þæt ic þe mínum eagum geseah.“\eva

\bvb TODO: Translation\evb\evg


\bvg\bva[14][821]%
Þá spræc Eue eft, \hld\ idesa scíenost, &
wífa wlitegost; \hld\ híe wæs ge·weorc godes, &
þeah héo þá on deofles cræft \hld\ be·droren wurde: &
\llap{„}Þu meaht hit mé witan, \hld\ wine mín Adam, &
wordum þínum; \hld\ hit þé þéah wyrs ne mæg &
on þinu\emph{m} hyge hreowan \hld\ þonne hit mé æt heortan deð.“ &
Hire þá Adam \hld\ and·swarode: &
\llap{„}Gif ic waldendes \hld\ willan cu̇ðe, &
hwæt ic his to hearm-sceare \hld\ habban sceolde, &
ne ge·sawe þú nó sniomor, \hld\ þeah mé on sǽ wadan &
hete heofones god \hld\ heonone nú þá, &
on flód faran, \hld\ nǽre hé firnum þæs deop, &
mere-stream þæs micel, \hld\ þæt his ó mín mód ge·tweode, &
ac ic tó þam grunde genge, \hld\ gif ic godes meahte &
willan ge·wyrcean. \hld\ Nis me on worulde niod &
æniges þegn-scipes, \hld\ nu ic mínes þeodnes hafa &
hyldo for·worhte, \hld\ þæt ic hie habban ne mæg. &
Ac wit þus baru ne magon \hld\ bú tú æt·somne &
wesan to wuhte. \hld\ Uton gán on þysne weald innan, &
on þisses holtes hleo.“ \hld\ Hwurfon hie bá twá, &
to·gengdon gnorngende \hld\ on þone grénan weald. &
Sǽton on·sundran, \hld\ bídan selfes ge·sceapu &
heofon-cyninges, \hld\ þá hie þá habban ne móston &
þe him ær for·geaf \hld\ æl-mihtig god. &
Þá hie heora líc-hǫman \hld\ leafum be·þeahton, &
weredon mid ðý wealde, \hld\ wǽda ne hæfdon; &
ac hie on ge·bed feollon \hld\ bú tú æt·somne &
morgena ge·hwilce, \hld\ bǽdon mihtigne &
þæt hie ne for·geate \hld\ god ælmihtig, &
and him ge·wísade \hld\ waldend se góda, &
hú hie on þam leohte forð \hld\ libban sceolden.\eva

\bvb TODO: Translation\evb\evg

\sectionline
